\part{Proof Theory}
\label{part:proof-theory}

\chapter[Classical and Intuitionistic Calculi]{Classical and Intuitionistic Sequent Calculi}
\label{ch:lk-lj}

\begin{convention}[Proof-theory formula metavariables]
Throughout this part, \(A,B,C\) range over formulas unless a section explicitly
declares a narrower local role.  The symbols \(\Gamma,\Delta,\Sigma,\Pi\)
retain their normative standing as formula contexts.  The simply typed lambda
calculus section explicitly redeclares \(A,B\) as type metavariables and
\(M,N\) as term metavariables only within that section.
\end{convention}

The sequent calculi \(\LK\) and \(\LJ\), and the cut-elimination
theorem of \cref{ch:cut-elimination}, originate with Gentzen
\cite{Gentzen1934a,Gentzen1935b}.

Structural proof theory studies derivations as formal objects.  The
active calculus, its context discipline, and its conversion relation must be
fixed before a claim about cut elimination or normalization is meaningful.
The presentation below follows the standard separation between sequent
calculus and natural deduction in
\cite{TroelstraSchwichtenberg2000,NegriVonPlato2001}.

\section{Sequents and contexts}
\label{sec:lk-lj-01}

\begin{convention}[Contexts]
In this chapter a context is a finite multiset of formulas.  Commas denote
multiset union.  Exchange is therefore built into the representation of a
context; weakening and contraction are not.  If \(\Gamma\) is a context, then
\(\Gamma,A\) adds one occurrence of \(A\).
\end{convention}

\begin{definition}[Classical sequent]
A classical sequent is an expression
\[
  \Gamma\Rightarrow\Delta,
\]
where \(\Gamma\) and \(\Delta\) are finite multisets of formulas.  Its semantic
reading is
\[
  \bigwedge\Gamma\to\bigvee\Delta,
\]
with \(\bigwedge\varnothing=\topv\) and
\(\bigvee\varnothing=\bottom\).  This reading states the intended semantics; a
sequent is itself a syntactic object, not an object-language conditional.
\end{definition}

\begin{definition}[Intuitionistic sequent]
An intuitionistic sequent has the form
\[
  \Gamma\Rightarrow A,
\]
with exactly one succedent formula.  The restriction to one succedent is part
of the definition of the calculus \(\LJ\), not a typographic abbreviation for
a classical sequent.
\end{definition}

